\begin{problem}
    Let $GL(n, \mathbf{R}) \subset \mathbf{R}^{n \times n}$ denote the general linear group of invertible $n \times n$ real matrices. This set is an open subset of the vector space $\mathbf{R}^{n \times n}$. Consider the inversion map $g: GL(n, \mathbf{R}) \to GL(n, \mathbf{R})$ defined by$$g(A) = A^{-1}.$$Prove that $g$ is  differentiable at any $A \in GL(n, \mathbf{R})$, and show that its total derivative $g'(A): \mathbf{R}^{n \times n} \to \mathbf{R}^{n \times n}$ is the linear map given by$$g'(A)(H) = -A^{-1} H A^{-1}.$$\\
 Hints:Algebraic Identity: For a small perturbation matrix $H$, start by factoring $A$ out of the expression $(A+H)^{-1}$ to relate it to the identity matrix:$$(A+H)^{-1} = (A(I + A^{-1}H))^{-1} = (I + A^{-1}H)^{-1} A^{-1}.$$ The Geometric Series: Recall that for a matrix $M$ with $\|M\| < 1$, the following power series (the Neumann series) converges:$$(I + M)^{-1} = I - M + M^2 - M^3 + \dots$$Apply this to $M = A^{-1}H$ for $H$ sufficiently small. Isolate the Linear Term: Expand $g(A+H) - g(A)$ using the series. Identify the term that is linear in $H$. This is your candidate for $g'(A)(H)$.
\end{problem}

\begin{problem}
    Let $f : \mathbf{R}^2 \to \mathbf{R}$ be defined by
$$
f(x,y)
:=
\frac{xy^3}{x^2+y^2}
\quad \text{when } (x,y)\ne(0,0),
\qquad
f(0,0):=0.
$$
\begin{itemize}
\item[(a)] Find $f_x(0,0)$ and $f_y(0,0)$.
 \item[(b)] Prove that $f$ is differentiable everywhere and that $\nabla f$ is continuous on $\mathbf{R}^2$, i.e.\ $f\in C^1(\mathbf{R}^2)$. Hint: You can use polar coordinate 
$x=r \cos(\theta), y=r \sin (\theta)$ when you try to show that  $f_x \to 0$ and $f_y \to 0$ as
$(x,y)\to 0$. Also $(x,y)\to 0$ implies $r\to 0$
\item[(c)] Show that the mixed partial derivatives $f_{xy}(0,0)$ and $f_{yx}(0,0)$ both exist.
\item[(d)] Compute $f_{xy}(0,0)$ and $f_{yx}(0,0)$ and show that they are \emph{not} equal.
\item[(e)] Identify precisely which hypothesis of Clairaut's theorem fails here.
\end{itemize}
\end{problem}

\begin{problem}
    Let $E\subset\mathbf{R}^2$ be open, and let $f\in C^2(E)$.
Fix $(x_0,y_0)\in E$.  For sufficiently small $h,k$ such that the rectangle
\[
R_{h,k}:=[x_0,x_0+h]\times [y_0,y_0+k]\subset E,
\]
define the second finite difference
\[
\Delta_{h,k} f
:=
f(x_0+h,y_0+k)-f(x_0+h,y_0)-f(x_0,y_0+k)+f(x_0,y_0).
\]
\begin{itemize}
\item[(a)] Prove the identity
\[
\Delta_{h,k} f
=
\int_{x_0}^{x_0+h}\int_{y_0}^{y_0+k} f_{xy}(x,y)\,dy\,dx
=
\int_{y_0}^{y_0+k}\int_{x_0}^{x_0+h} f_{yx}(x,y)\,dx\,dy.
\]
(You may use the one-variable Fundamental Theorem of Calculus twice.)
\item[(b)] Deduce that
\[
\frac{\Delta_{h,k} f}{hk}\to f_{xy}(x_0,y_0)
\quad\text{and}\quad
\frac{\Delta_{h,k} f}{hk}\to f_{yx}(x_0,y_0)
\qquad\text{as }(h,k)\to(0,0),
\]
and conclude $f_{xy}(x_0,y_0)=f_{yx}(x_0,y_0)$. Hint:  If $F$ is continuous on $[x_0,x_0+h]\times [y_0,y_0+k]$, then there exists
$(\bar x,\bar y)\in [x_0,x_0+h]\times [y_0,y_0+k]$ such that
\[
\int_{x_0}^{x_0+h}\int_{y_0}^{y_0+k} F(x,y)\,dy\,dx
= hk\,F(\bar x,\bar y).
\]
Similarly, there exists $(\tilde x,\tilde y)\in [x_0,x_0+h]\times [y_0,y_0+k]$ such that
\[
\int_{y_0}^{y_0+k}\int_{x_0}^{x_0+h} F(x,y)\,dx\,dy
= hk\,F(\tilde x,\tilde y).
\]

\end{itemize}

\medskip
\noindent\emph{Hint:} For (b), use continuity of $f_{xy}$ (resp.\ $f_{yx}$) to compare the double integral with $hk\,f_{xy}(x_0,y_0)$ (resp.\ $hk\,f_{yx}(x_0,y_0)$).
\end{problem}

\begin{proof}[(a)]
    By applying Fundamental Theorem of Calculus, we have 
    \begin{align*}
        \int _{x_0}^{x_0 + h} \int _{y_0}^{y_0 + k} f_{xy}(x, y) \mathrm{d} y \mathrm{d} x &= \int _{x_0}^{x_0 + h} \left. f_x(x, y) \right]_{y_0}^{y_0 + k} \mathrm{d} x = \int _{x_0}^{x_0 + h} f_x (x, y_0 + k) - f_x(x, y_0) \mathrm{d} x \\
        &= \left[ f(x, y_0 + k) - f(x, y_0) \right]_{x_0}^{x_0 + h} \\
        &= \left( f(x_0 + h, y_0 + k) - f(x_0 + h, y_0) \right) - \left( f(x_0, y_0 + k) - f(x_0, y_0) \right) \\
        &= f(x_0 + h, y_0 + k) - f(x_0 + h, y_0) - f(x_0, y_0 + k) + f(x_0, y_0) \\
        &= \Delta _{h, k} f      
    \end{align*}    
    and 
    \begin{align*}
        \int _{y_0}^{y_0 + k} \int _{x_0}^{x_0 + h} f_{yx}(x, y) \mathrm{d} x \mathrm{d} y &= \int _{y_0}^{y_0 + k} \left. f_y(x, y) \right]_{x_0}^{x_0 + h} \mathrm{d} y = \int _{y_0}^{y_0 + k} f_y (x_0 + h, y) - f_y(x_0, y) \mathrm{d} y \\
        &= \left[ f(x_0 + h, y) - f(x_0, y) \right]_{y_0}^{y_0 + k} \\
        &= f(x_0 + h, y_0 + k) - f(x_0 + h, y_0) - f(x_0, y_0 + k) + f(x_0, y_0) \\
        &= \Delta _{h, k} f.      
    \end{align*}
    Thus,
    \[
        \Delta_{h,k} f
=
\int_{x_0}^{x_0+h}\int_{y_0}^{y_0+k} f_{xy}(x,y)\,dy\,dx
=
\int_{y_0}^{y_0+k}\int_{x_0}^{x_0+h} f_{yx}(x,y)\,dx\,dy.
    \]
\end{proof}

\begin{proof}[(b)]
    We know \(R_{h, k}\) is closed and bounded on \(\mathbb{R} ^2\), so \(R_{h, k}\) is compact by Heine Borel's theorem. Also, for sufficiently small \(h, k\), we know \(R_{h, k} \subseteq E\) since \(E\) is open, so in this case \(f_{xy}\) is continuous on \(R_{h, k}\). Now since \(f_{xy}\) is continuous on \(R_{h, k}\) and \(R_{h, k}\) is compact, so \(\exists v_{\text{max}}, v_{\text{min}} \in R_{h, k}\) s.t. 
    \[
        f_{xy}(v_{\text{max}}) = \max _{v \in R_{h, k}} f_{xy}(v), \quad f_{xy} (v_{\text{min}}) = \min_{v \in R_{h, k}} f_{xy}(v). 
    \]          
    Now let \(f_{xy}(v_{\text{max}}) = M\) and \(f_{xy} (v_{\text{min}}) = m\), then for all \(v \in R_{h, k}\), \(m \le f_{xy}(v) \le M\). Thus, we have 
    \[
        hk \cdot m \le \int _{x_0}^{x_0 + h} \int _{y_0}^{y_0 + k} f_{xy}(x, y) \mathrm{d} y \mathrm{d} x \le hk \cdot M.  
    \]    
    Note that \(R_{h, k}\) is connected and \(f_{xy}\) is continuous, so \(f_{xy} (R_{h, k})\) is a connected subset of \(\mathbb{R} \), i.e. \(f_{xy} (R_{h, k}) = [m, M]\), so there exists \(\left( \overline{x} , \overline{y}  \right) \in R_{h, k} \) s.t. 
    \[
        f_{xy} \left( \overline{x} , \overline{y}  \right) = \frac{1}{hk} \int _{x_0}^{x_0 + h} \int _{y_0}^{y_0 + k} f_{xy} (x, y) \mathrm{d} y \mathrm{d} x = \frac{1}{hk} \Delta _{h, k} f.   
    \]      
    Now if \((h, k) \to (0, 0)\), then \((\overline{x} , \overline{y} ) \to (x_0, y_0)\), and since \(f_{xy}\) is continuous on \(R_{h, k}\), we have \(f_{xy} \left( \overline{x} , \overline{y}  \right) \to f_{xy}(x_0, y_0) \). Thus, we can conclude 
    \[
        \frac{1}{hk} \Delta _{h, k} f \to f_{xy} (x_0, y_0) \text{ as } (h, k) \to (0, 0). 
    \]
    By using same arguments but replacing \(f_{xy}\) by \(f_{yx}\) and use 
    \[
        \Delta_{h,k} f = \int_{y_0}^{y_0+k}\int_{x_0}^{x_0+h} f_{yx}(x,y)\,dx\,dy,
    \]       
    we can deduce 
    \[
        \frac{1}{hk} \Delta _{h, k} f \to f_{yx} (x_0, y_0) \text{ as } (h, k) \to (0, 0).
    \]
    Finally, we can conclude 
    \[
        f_{xy}(x_0, y_0) = f_{yx} (x_0, y_0).
    \]
\end{proof}

\begin{problem}
    Let $f$ and $g$ be functions from $\mathbf{R}^n$ to $\mathbf{R}^m$. Assume that $f$ is
differentiable at $c$, that $f(c)=0$, and that $g$ is continuous at $c$.
Let
\[
h(x)=g(x)\cdot f(x).
\]
Prove that $h$ is differentiable at $c$ and that
\[
h'(c)(u)=g(c)\cdot\bigl(f'(c)(u)\bigr)
\qquad \text{if } u\in \mathbf{R}^n.
\] Note that $g(x)\cdot f(x)$ is the inner product of two vectors in $\mathbf{R}^m$ and $h$ is a function from
$\mathbf{R}^n$ to $\mathbf{R}$.
\end{problem}

\begin{problem}
    Let $f$ and $g$ be real-valued functions defined on $\mathbf{R}$ with continuous second
derivatives $f''$ and $g''$. Define
\[
F(x,y)=f\bigl(x+g(y)\bigr)
\qquad \text{for each }(x,y)\in \mathbf{R}^2.
\]
Find formulas for all first- and second-order partial derivatives of $F$ in terms of the
derivatives of $f$ and $g$. Verify the relation
\[
F_x\,F_{xy}=F_y\,F_{xx}.
\]
\end{problem}